\section{Using with \handbookMistborn{}}
\label{sec:otherRules}

This set of rules is meant to expand upon the content of \handbookMistborn{}, not replace the rules defined there. The following is still true for Hemalurgic spikes.

\subsection{Implanting and Removing Spikes}

This set of rules defines how a Hemalurgist can create new Spikes. That doesn't change the rules of implanting the Spikes described in \handbookMistborn{}.

For more details, see section "Implanting and Removing Hemalurgic Spikes" (\handbookMistborn{}, page 289.).

\subsubsection{Implant Spike (\iconActionTriple[8pt]{})}

To implant a created Spike, make a \skillCheck{20}{\skillMedicine} targeting yourself or a willing allay. On a success, you place the Spike in the correct binding point, and its effects take place. On a failure, the target takes 1d10 vital damage, and the spike is not implanted.

\subsubsection{Remove Spike (\iconAction[8pt]{})}

You can attempt to remove a Spike from you or a willing ally you can touch. To do so, make a \skillCheck{10}{\skillMedicine}. On a failure, the target takes 1d4 vital damage and the spike remains in place.

\subsubsection{Removing an Enemy’s Spike}

\handbookMistborn{} describes multiple methods of removing Hemalurgic Spikes from your enemies. All of them still apply when using this set of homebrew rules. Work with your GM to determine which method is best suited for your current situation.



\subsection{Hemalurgic Spike Effects}

While this set of rules defines some traits and effects a Hemalurgic Spike obtains during its creation, all effects defined in \handbookMistborn{} still take place. This book expands them, but is not meant to replace them.

For more details, see section "Hemalurgic Spike Effects" (\handbookMistborn{}, page 290.).

\subsubsection{Spiritweb Disruption}

By default, a PC can have up to 3 Hemalurgic Spikes implanted at the same time. When the 4th one is implanted, the PC's spiritweb becomes disrupted and they succumb to the full control of the Shard. At that point, that character becomes an NPC and the player must create a new PC to continue playing.

However, when using this ruleset, a skilled Hemalurgist can create a Linchpin to increase that limit.

\subsubsection{Influence of Shards}

Whispers of Ruin trait on Hemalurgic Spikes was designed to be a clearly defined set of conditions and test DCs to represent the influence of the Shards and help the GM to define and trigger those moments.

This doesn't mean the GM cannot initiate the influence on their own, if the story demands it. The GM can also set their own DC for the tests. The rules in this book are meant to help the GM, not constrict them, and, as always, the GM has the final say.
